\section{Metodología de Evaluación}

La objetividad de la evaluación y la verificación del cumplimiento de los KPIs se fundamentan en la transición de registros analógicos manuales hacia fuentes de información digitalizadas y centralizadas. Para contrastar el estado inicial (Línea Base) con los resultados post-implementación, se ha estructurado una metodología basada en tres ejes fundamentales: fuentes de datos, participantes involucrados y un cronograma de medición.

\subsection{Fuentes de Datos}
La recolección de métricas se realiza de manera objetiva a través de las siguientes herramientas del nuevo ecosistema digital:
\begin{itemize}
    \item \textbf{Exactitud de Inventario:} Se contrasta el reporte comparativo de cierre de stock emitido por el sistema POS de Odoo contra las auditorías físicas periódicas de control.
    \item \textbf{Tiempo de Respuesta (Provincias):} Se extraen datos directamente de los logs de atención y capturas de pantalla integradas en el sistema de gestión.
    \item \textbf{Presencia Digital y Omnicanalidad:} Se analizan las métricas nativas de los paneles de \textit{Meta Business Suite} (para Facebook e Instagram) y \textit{WhatsApp Business}, verificando la operatividad de los enlaces corporativos.
    \item \textbf{Toma de Decisiones y Reportabilidad:} Se evalúa la disponibilidad y uso del Reporte Semanal Automático de "productos estrella", generado en tiempo real por los dashboards del software de gestión.
\end{itemize}

\subsection{Participantes de la Evaluación}
La validación de usabilidad y el control de los indicadores involucran activamente a los actores clave del flujo de comercialización:
\begin{itemize}
    \item \textbf{Fuerza de Ventas y Almacén (equipo de 5 colaboradores):} Como usuarios directos del POS e inventario \textit{cloud}, se evalúa su nivel de adopción tecnológica y el reemplazo definitivo de los cuadernos físicos. El Encargado de Ventas/Almacén reporta directamente sobre la precisión del registro de salidas.
    \item \textbf{Responsable de Marketing:} Monitorea el alcance de las campañas, el tráfico hacia el nuevo catálogo digital y la sincronización del stock visible en línea.
    \item \textbf{Gerencia General:} Usuario estratégico que valida los resultados consolidados a través de los dashboards, aportando \textit{feedback} fundamental sobre la utilidad del sistema para optimizar compras y planificación operativa.
\end{itemize}

\subsection{Periodo y Frecuencia de Medición}
El plan de evaluación contempla un horizonte estructurado para garantizar la fiabilidad de los resultados medidos:
\begin{itemize}
    \item \textbf{Línea Base:} Definida formalmente en abril de 2026, consolidando los datos de la situación analógica previa a la automatización.
    \item \textbf{Evaluación Inicial (Marcha Blanca):} El periodo de estabilización operativa, puesta en marcha del POS y levantamiento de inventario abarca de abril a junio de 2026.
    \item \textbf{Horizonte Completo y Frecuencia:} El proyecto global se enmarca en un horizonte de 6 meses. La medición de KPIs se realiza de forma escalonada: \textbf{diaria/semanal} para stock crítico y atención al cliente; \textbf{mensual} para tasas de error y efectividad de pauta; y \textbf{semestral} para la auditoría final de inventario y rentabilidad global.
\end{itemize}
